% Part: applied-modal-logics
% Chapter: epistemic-logic
% Section: bisimulations

\documentclass[../../../include/open-logic-section]{subfiles}

\begin{document}

\olfileid{aml}{el}{bsd}

\olsection{દ્વિઅનુકરણો}

આપણી જ્ઞાનસંબંધી ભાષાની અભિવ્યક્તિ-શક્તિ વિશેનો એક બાકી પ્રશ્ન
નિદર્શો અને તેમાં સત્ય હોય એવાં સૂત્રો વચ્ચેના સંબંધને લગે છે.
ફ્રેમ-અનુરૂપતાનાં પરિણામોમાં આપણે જોયું છે કે અમુક સૂત્રો
ફ્રેમમાં પ્રામાણ્ય હોય, તો એ ફ્રેમ અમુક ગુણધર્મો ધરાવે છે.
પણ શું આપણી મોડલ ભાષા, દાખલા તરીકે, પોતાને જ જતું સ્વવાચક તીર
ધરાવતા જગત અને એક પછી બીજા જગતમાં જતી અનંત શૃંખલા વચ્ચે ભેદ
કરી શકે? એટલે કે, આ બેમાંથી માત્ર એક જ જગતમાં સત્ય હોય એવું
કોઈ સૂત્ર~$A$ છે?

\emph{દ્વિઅનુકરણ} (બાઇસિમ્યુલેશન) એવો સંબંધ છે જે બે
સંબંધાત્મક નિદર્શો વચ્ચે વ્યાખ્યાયિત કરી શકાય છે અને દર્શાવે
છે કે એમની રચના અસરકારક રીતે સમાન છે. આપણે જોઈશું કે તે
આપણી જ્ઞાનસંબંધી ભાષા વડે ઓળખી શકાય એવી નિદર્શોની એક પ્રકારની
સમતુલ્યતા પકડે છે.

\begin{defn}[દ્વિઅનુકરણ]
$M_1 = \tuple{W_1, R_1, V_1}$ અને $M_2 = \tuple{W_2, R_2, V_2}$
એ બે સંબંધાત્મક નિદર્શો હોય. વળી $\mathcal{R} \subseteq W_1 \times W_2$
દ્વિસ્થાની સંબંધ હોય. દરેક $\tuple{w_1, w_2} \in \mathcal{R}$
માટે નીચેની શરતો પૂરી થાય ત્યારે $\mathcal{R}$ને
\emph{દ્વિઅનુકરણ} કહીએ છીએ:
\sourcecorrection{OLMD-079}{મૂળની બીજી અને ત્રીજી શરતમાં
કર્તાઓનો ગણ $A$ લખાયો છે; આ પ્રકરણમાં તે ગણ $G$ તરીકે
વ્યાખ્યાયિત થયો છે. બંને શરતમાં $G$ મૂક્યો છે.}

\begin{enumerate}
  \item બધા !!{propositional variable}s~$p$ માટે
  $w_1 \in V_1(p)$ ત્યારે અને માત્ર ત્યારે જ જ્યારે $w_2 \in V_2(p)$.
  \item બધા કર્તાઓ $a \in G$ અને જગતો $v_1 \in W_1$ માટે,
  $R_{1_a} w_1 v_1$ હોય તો કોઈ $v_2 \in W_2$ એવું હોય કે
  $R_{2_a} w_2 v_2$ અને $\tuple{v_1, v_2} \in
  \mathcal{R}$.
  \item બધા કર્તાઓ $a \in G$ અને જગતો $v_2 \in W_2$ માટે,
  $R_{2_a} w_2 v_2$ હોય તો કોઈ $v_1 \in W_1$ એવું હોય કે
  $R_{1_a} w_1 v_1$ અને $\tuple{v_1, v_2} \in
  \mathcal{R}$.
\end{enumerate}

$M_1$ અને $M_2$ વચ્ચેના દ્વિઅનુકરણમાં જગતો $w_1$ અને $w_2$
જોડાયેલા હોય, તો $\tuple{M_1, w_1} \leftrightarroweq
\tuple{M_2, w_2}$ પણ લખી શકીએ છીએ અને $\tuple{M_1, w_1}$ તથા
$\tuple{M_2, w_2}$ને \emph{દ્વિઅનુકરણ-સમાન} કહીએ છીએ.
\end{defn}

દ્વિઅનુકરણની જુદી શરતો જુદી બાબતો સુનિશ્ચિત કરે છે.
પહેલી શરત બધા !!{propositional variable}s વિશે સહમતિ
સુનિશ્ચિત કરે છે; તેથી દ્વિઅનુકરણ-સમાન જગતોમાં સમાન
મોડલ-મુક્ત સૂત્રો સત્ય હશે. બીજી અને ત્રીજી શરતોને અનુક્રમે
``આગળ'' અને ``પાછળ''ની શરત કહેવાય છે; તે સુલભતા સંબંધોની
રચના સમાન હોવાની ખાતરી આપે છે.

\begin{thm}
જો $\tuple{M_1, w_1} \leftrightarroweq \tuple{M_2, w_2}$, તો દરેક
!!{formula}~$!A$ માટે $\mSat{M_1}{!A}[w_1]$ ત્યારે અને માત્ર
ત્યારે જ જ્યારે $\mSat{M_2}{!A}[w_2]$.
\end{thm}

\begin{figure}
  \begin{center}
    \begin{tikzpicture}[modal]
      \node[world] (w1) {$w_1$}; 
      \node[world] (w2) [above left=of w1]{$w_2$}; 
      \node[world] (w3) [above right=of w1] {$w_3$};
      \draw[<->] (w1) to node {$a$} (w2);
      \draw[->,loop below] (w1) to node {$a$} (w1);
      \draw[<->] (w1) to [swap] node {$a$} (w3);
      \draw[->,loop above] (w2) to node {$a$} (w2) ;
      \draw[->,loop above] (w3) to node {$a$} (w3) ;
      
      \node[world](v1)[right of=w1, xshift=1.5in]{$v_1$};
      \node[world](v2)[above of=v1]{$v_2$};
      \draw[<->] (v1) to node {$a$} (v2);
      \draw[->,loop below] (v1) to node {$a$} (v1);
      \draw[->,loop above] (v2) to node {$a$} (v2) ;

      \draw[-,dotted] (w1) to (v1) ;
      \draw[-,dotted,bend left=45] (w2) to (v2) ;
      \draw[-,dotted,bend left=30] (w3) to (v2) ;
    \end{tikzpicture}
  \end{center}
  \caption{બે દ્વિઅનુકરણ-સમાન નિદર્શો.}
  \ollabel{fig:bisimilar}
\end{figure}

\olref{fig:bisimilar}માં દર્શાવેલા બે નિદર્શો એકસરખા નથી,
છતાં જગતો $w_1$ અને~$v_1$ને જોડતું દ્વિઅનુકરણ છે. આ
દ્વિઅનુકરણ $w_2$ અને $w_3$ બંનેને~$v_2$ સાથે પણ જોડે છે:
આપણી મોડલ ભાષામાં વ્યક્ત થઈ શકે એવી કોઈ બાબત ખરેખર એમની
વચ્ચે ભેદ કરી શકતી નથી. જોકે $w_2$ અને~$w_3$માં જુદાં
!!{propositional variable}s સત્ય હોત, તો સ્થિતિ જુદી હોત.

\end{document}
